\section{Linear Transformations}
 (HERE COMES LOTS OF INTRO TEXT ON LINEAR TRANSFORMATIONS WHICH I WILL WRITE LATER)

Linear transformations have a very nice property which makes applying them on different vectors extremely efficient. Consider the linear transformation $T:\Rs[2]\to\Rs[2]$ which is a composition of scaling by a factor of $2$ in the $y$-diretion followed by a rotation of $\ang{90}$ anti-clockwise around the origin. Now let's say that we wish to calculate how $T$ transforms the vector
$$\vec{a}=\colvec{3;1}.$$

In principle, we could simply apply the operations one after the other on $\vec{a}$: scaling by $2$ in the $y$-direction yields the vector $$\vec{b}=\colvec{3;2},$$ and the the rotation gives the final vector
\begin{equation}
	T(\vec{a}) = \colvec{-2;3}.
	\label{eq:simple_LT_example_single_2d_vec}
\end{equation}

Now, what happens if we wish to calculate a more complicated linear transformation, which is both composed out of many more basic transformations - and also acts on higher-dimensional vectors? Performing each transformation step might not be so straight-forward, and is definitely not an efficient process. In fact, even if we're dealing with a simple example just like the one we just saw, we would still need to apply it every time we want to calculate how a new vector is transformed. It turns out that thanks to the simplicity of linear transformations \footnote{more correctly: to the strict rules they enforce} we only need to perform a few simple calculations per transformation to practically get any possible application of it on any vector!

To see how this is the case, let's go back to our simple $2$-dimensional example. First, recall that we can express any vector in terms of the standard basis vectors (REF), which in this case gives $$\vec{a}=\colvec{3;1}=3\colvec{1;0}+1\colvec{0;1}$$. We can then apply the defining properties of linear transformations (REF) on $\vec{a}$:

\begin{equation}
	\begin{aligned}
		T\colvec{3;1} & = T\left(3\colvec{1;0}+1\colvec{0;1}\right)                 \\
		              & = T\left(3\colvec{1;0}\right) + T\left(1\colvec{0;1}\right) \\
		              & = 3T\colvec{1;0} + 1T\colvec{0;1}                           \\
		              & = 3T\left(\eb{1}\right) + T\left(\eb{2}\right).
	\end{aligned}
	\label{eq:}
\end{equation}

Now, applying $T$ to the two standard basis vectors $\eb{1}$ and $\eb{2}$ we get
\begin{equation}
	\begin{aligned}
		T\left(\eb{1}\right) & = \colvec{0;1},  \\
		T\left(\eb{2}\right) & = \colvec{-2;0}, \\
	\end{aligned}
	\label{eq:}
\end{equation}
and then
\begin{align*}
	3T\left(\eb{1}\right) + T\left(\eb{2}\right) & = 3\colvec{0;1} + \colvec{-2;0} \\
	                                             & = \colvec{0;3} + \colvec{-2;0}  \\
	                                             & = \colvec{-2;3},
\end{align*}
which is exactly what we got when we calculated $T$'s effect directly on $\vec{a}$ (cf. \autoref{eq:simple_LT_example_single_2d_vec}).
